\documentclass[12pt]{article}
\usepackage[utf8]{inputenc}
\usepackage[T1]{fontenc}
\usepackage[spanish]{babel}
\usepackage{geometry}
\usepackage{amsmath,amssymb}
\usepackage{enumitem}

\geometry{margin=2.5cm}

\begin{document}

\begin{center}
{\Large Transcripción cruda}\\[2mm]
{\large 18May26}
\end{center}

\section*{Foto 1}

Fecha para el primer examen parcial: jueves de semana 5.

1er lista de ejercicios: capítulo II de Rudin, del 1 al 19.

\textit{Teorema.}
\begin{enumerate}[label=\alph*)]
    \item Intersección finita de conjuntos abiertos es un conjunto abierto.
    \item Unión finita de conjuntos cerrados es un conjunto cerrado.
\end{enumerate}

\textit{Ejemplo.}
\begin{enumerate}[label=\alph*)]
    \item Intersección infinita de abiertos no es un conjunto abierto.
    \item Unión infinita de conjuntos cerrados no es cerrada.
\end{enumerate}

\[
\forall n\in\mathbb N,\quad G_n=\left(-\frac1n,\frac1n\right)
\]
es abierto en $\mathbb R$ y
\[
\bigcap_{n=1}^{\infty}G_n=\{0\}.
\]

\[
G_n^c=\left(-\infty,-\frac1n\right]\cup\left[\frac1n,\infty\right),
\qquad
\bigcup_{n=1}^{\infty}G_n^c=(-\infty,0)\cup(0,\infty),
\]
que es abierto pero no cerrado.

\section*{Fotos 2--5}

Definiciones sobre interior, cerradura y conjunto derivado.

Sea $(X,d)$ un espacio métrico y sea $E\subset X$.

\[
E'=\{x\in X\mid x \text{ es punto límite de }E\}.
\]

En la definición anterior, la cerradura de $E$ es
\[
\overline E:=E\cup E'.
\]

Consecuencia de la proposición:
\[
\bigcup_{A\in\mathcal D}A=E^\circ,
\]
donde $\mathcal D=\{A\subset E\mid A \text{ es abierto}\}$.

\section*{Fotos 6--10}

Ejemplo:
\[
E=\{(x,y)\in\mathbb R^2\mid x^2+y^2\le 1,\ y>0\}.
\]
El conjunto $E$ no es cerrado ni es abierto.

\[
E^\circ=\{(x,y)\in\mathbb R^2\mid x^2+y^2<1,\ y>0\}.
\]

\[
\overline E=\{(x,y)\in\mathbb R^2\mid x^2+y^2\le 1,\ y\ge 0\}.
\]

\[
E'=\overline E.
\]

\section*{Fotos 11--14}

Teorema: Sea $(X,d)$, sea $Y\subset X$ y $E\subset Y$.
Entonces $E$ es abierto relativo en $Y$ si y solo si
\[
E=Y\cap G
\]
con $G$ un abierto en $X$.

\section*{Fotos 15--19}

Sea $E\subset X$.
Si $y\in E$, entonces $y$ no es cota superior de $E$.
Si $y\notin E$ y $y$ es cota superior de $E$, entonces para todo $\varepsilon>0$ existe $x\in E$ tal que
\[
y-\varepsilon<x\le y.
\]
Si además $y\notin E$, entonces
\[
y-\varepsilon<x<y.
\]
Por tanto
\[
N_\varepsilon(y)\cap(E\setminus\{y\})\neq\varnothing,
\]
y entonces $y\in E'$.

Algo similar ocurre con el infinito de un subconjunto en $\mathbb R$ no vacío y acotado inferiormente.

\section*{Fotos 20--23}

\textit{Definición.} Sea $(X,d)$ un espacio métrico y sean $E,Y\subset X$ tales que $E\subset Y\subset X$.
Diremos que $E$ es abierto relativo en $Y$ si para todo $y\in E$ existe $r>0$ tal que
\[
N_r(y)\cap Y\subset E.
\]

\textit{Teorema.} Sea $(X,d)$ un espacio métrico y sea $Y\subset X$.
Entonces $E$ es abierto relativo en $Y$ si y solo si $E=Y\cap G$ con $G$ abierto en $X$.

\section*{Fotos 24--25}

\textit{Teorema.} Si $(X,d)$ es un espacio métrico y $K\subset X$ es compacto, entonces $K$ es cerrado.

Demostración: probaremos que $K^c$ es abierto.

Sea $p\in K^c$. Para cada $q\in K$, sea
\[
r_q=\frac12 d(p,q)>0.
\]
Entonces
\[
N_{r_q}(q)\cap N_{r_q}(p)=\varnothing.
\]
Como $K$ es compacto en $(X,d)$, existe una subcubierta finita de $\{N_{r_q}(q)\}_{q\in K}$.
Si $q_1,\dots,q_n\in K$, entonces
\[
K\subset N_{r_{q_1}}(q_1)\cup\cdots\cup N_{r_{q_n}}(q_n).
\]
Tomando
\[
r=\min\{r_{q_1},\dots,r_{q_n}\},
\]
se obtiene
\[
N_r(p)\cap K=\varnothing.
\]
Por tanto $K^c$ es abierto y $K$ es cerrado.

\end{document}
